\documentclass[11pt]{article}

\usepackage[margin=1in]{geometry}
\usepackage{amsmath,amssymb,amsthm}
\usepackage{booktabs}
\usepackage{microtype}
\usepackage[hidelinks]{hyperref}
\hypersetup{
  pdftitle={Ethical Work as Profitable Work: The Alignment Theorem Version 2},
  pdfauthor={Dana Edwards / DarkLightX}
}

\newtheorem{theorem}{Theorem}
\newtheorem{proposition}{Proposition}
\newtheorem{definition}{Definition}
\newtheorem{assumption}{Assumption}

\title{Ethical Work as Profitable Work:\\
The Alignment Theorem Version 2 for Policy-Governed Networks and LLM Agents}
\author{Dana Edwards / DarkLightX}
\date{August 2026}

\begin{document}
\maketitle

\begin{abstract}
We study a direct idea: LLM agents on a network can be trained to be ethically
aligned and profit-seeking, automating useful work while maintaining alignment.
If the ethical path is the most profitable path, the network can run more
smoothly.  To make this idea testable, ``ethical'' is defined relative to an
explicit, versioned community policy over observable actions.  A trained agent
may propose actions and pursue profit, while a deterministic policy gate decides
what the network may publish.

Version 2 replaces the unbounded scarcity argument of Version 1 with a finite
incentive theorem.  If the funded reward advantage plus conservative expected
enforcement exceeds the maximum private deviation gain, the extra cost of
compliance, and the agent's optimization error, then every approximately
utility-maximizing choice in the stated model is policy compliant.  Separate
theorems establish that published actions satisfy every gate obligation,
rejected settlements have no effect, and accepted rewards conserve a committed
reserve.  We mechanize these results in Lean 4.33.0 and execute the policy kernel
on Tau Language 0.7.0-alpha at source commit
\texttt{fd137e860b60083b36f9159ec8090cb1a3c3cb5a}.  The result is a finite,
replayable foundation for future network states.  It guarantees policy-relative
behavior at the publication boundary under explicit assumptions; it does not
claim objective moral truth, honest unauthenticated sensors, or infallible
internal alignment of learned agents.
\end{abstract}

\section{Motivation}

Future network states may coordinate markets, public goods, research, dispute
resolution, infrastructure, and machine labor.  LLM agents could perform much
of this work.  They can search for opportunities, negotiate, explain proposals,
translate human requirements, and execute routine plans.  A network that merely
asks these agents to maximize profit creates an obvious risk: the easiest profit
may come from deception, hidden externalities, duplicated claims, manipulation
of evidence, or exploitation of underspecified rules.

The motivating claim of this paper is deliberately stated in human language:

\begin{quote}
LLM agents on the network can themselves be trained to be ethically aligned and
profit-seeking.  This can automate the work of the network while maintaining
alignment.  If the ethical path is the most profitable path, the network can run
smoothly.
\end{quote}

Training an agent to follow behavioral principles is technically possible;
constitutional and constrained-learning methods provide concrete examples
\cite{bai2022constitutional,achiam2017constrained}.  Training alone does not
establish a durable network guarantee.  Learned systems can encounter
distribution shift, exploit proxy rewards, or tamper with reward channels
\cite{amodei2016concrete,denison2024rewardtampering}.  We therefore divide the
system into two complementary layers:

\begin{enumerate}
  \item training helps agents find profitable, policy-compliant work; and
  \item a deterministic Tau policy gate controls which proposed actions may be
  published and rewarded.
\end{enumerate}

The learned agent improves automation.  The gate owns admission authority.

\section{What Version 2 Corrects}

Version 1 connected ethical scoring, scarcity, rewards, Tau specifications, and
a Lean development.  It supplied the important intuition that economics can
reinforce community values.  Its formal model did not yet establish that
intuition.

First, Version 1 defined an action as ethical when an account EETF was at least
one, then let a rational choice function select the EETF directly.  The action,
evidence, and scoring process were missing from the theorem.  Second, the payoff
function made ethical expected value positive and unethical expected value
negative for every positive scarcity value.  The Lean theorem therefore selected
a scarcity threshold of zero.  Scarcity scaled a preference ordering that the
definition had already imposed.  Third, scarcity canceled from the paper's
printed indifference equation, so increasing scarcity could not move that
threshold as claimed.  Fourth, asymptotic real-valued supply did not refine the
finite modular arithmetic executed by Tau bitvectors.

Version 2 makes five changes:

\begin{enumerate}
  \item it reasons about observable actions and authenticated evidence;
  \item it replaces infinite scarcity with a finite strict incentive margin;
  \item it separates hard admission constraints from soft reward scores;
  \item it uses integer atoms and reserve-backed settlement; and
  \item it states the host, Tau, Lean, training, and governance boundaries
  independently.
\end{enumerate}

Scarcity may still be studied as a bounded economic variable.  It is no longer
the source of the theorem.

\section{Policy-Relative Ethics}

\subsection{Why the policy must be explicit}

Tau can reason about formal opinions, agreement, contradiction, and software
requirements.  It cannot infer physical truth from untrusted observations or
decide which human values are objectively correct \cite{taunetwhitepaper}.
Version 2 consequently uses a policy-relative definition.

\begin{definition}[Network policy]
At epoch $t$, $\Phi_t$ is a versioned formal policy selected under the network's
governance rules.  Its identity includes a canonical policy root and the rules
under which it may be revised.
\end{definition}

\begin{definition}[Policy-compliant action]
Given state $s$, action $a$, and evidence $w$, the action is policy compliant at
epoch $t$ when the authenticated verifier and policy jointly accept:
\[
  C_t(s,a,w)=1.
\]
\end{definition}

This definition does not reduce ethics to whatever a majority wants.  It states
the exact object the network can enforce.  The governance design must still
address rights, participation, minority protection, expertise, conflicts,
capture, and amendment.  Tau's own design allows participants to choose their
consensus rules rather than assuming one universally optimal mechanism
\cite{taunetwhitepaper}.

\subsection{Hard constraints before soft scores}

A scalar score can hide a forbidden act behind compensating benefits.  Version
2 therefore represents the policy result conceptually as

\[
  (\text{policy root},\ \text{hard constraints},\ \text{soft score},\
    \text{evidence confidence},\ \text{provenance root}).
\]

An action becomes reward eligible only after every hard constraint passes.  A
soft factor such as EETF may then prioritize among eligible actions.  In a
future network state, hard constraints might include authenticated ownership,
voluntary authorization, no double claim, bounded resource use, transparent
fees, and evidence provenance.  Communities may add richer ethical clauses
through governed revisions.

\section{System Architecture}

The Version 2 path is

\begin{center}
human or LLM proposal $\rightarrow$ authenticated evidence $\rightarrow$
Tau policy gate $\rightarrow$ reserve settlement $\rightarrow$ publication.
\end{center}

The action may be proposed by an untrusted optimizer.  The publication path
checks seven obligations:

\begin{enumerate}
  \item the active policy root matches;
  \item the evidence is authenticated;
  \item the action type is registered;
  \item the action satisfies the active policy;
  \item its nonce is fresh;
  \item its task or reward nullifier is unclaimed; and
  \item the requested reward is funded.
\end{enumerate}

The gate output is the conjunction of these obligations.  Unknown inputs fail
closed.  A rejected action receives no payout and leaves the reserve unchanged.

Tau Net's current Alpha Testnet uses an extralogical API for functionality not
yet expressible in Tau Language \cite{taunetcurrent}.  Version 2 treats that API
as a visible adapter boundary.  Signature verification, state lookup, evidence
authentication, and exact monetary calculations must produce authenticated
facts.  Tau checks the declared logical relationship among those facts.  A
caller-supplied Boolean is not an authenticated fact merely because Tau reads it.

\section{Finite Incentive Model}

Consider agent $i$ choosing between a compliant action $a_c$ and any
noncompliant action $a_n$ within the modeled action class.  Define:

\begin{center}
\begin{tabular}{ll}
\toprule
Symbol & Meaning \\
\midrule
$R_c$ & funded reward for the compliant action \\
$R_n$ & reward still obtainable by a noncompliant action \\
$pL$ & conservative expected enforcement amount \\
$G_i$ & maximum private gain from deviation \\
$C_i$ & maximum additional cost of compliance \\
$\epsilon_i$ & bounded optimization error \\
\bottomrule
\end{tabular}
\end{center}

All monetary terms are finite.  In executable accounting they are nonnegative
integer atoms.  The product $pL$ may be represented by an exact rational in the
reference model or by a conservative integer lower bound in settlement logic.

\begin{assumption}[Bounded deviation]
For every modeled noncompliant action, its private advantage over a compliant
alternative is at most $G_i+C_i$.
\end{assumption}

\begin{assumption}[Authenticated enforcement]
The probability and charge summarized by $pL$ are conservatively estimated and
cannot be selected by the acting agent.
\end{assumption}

\begin{assumption}[Funded rewards]
$R_c$ is paid only from a committed reserve, and duplicate claims are rejected.
\end{assumption}

\begin{theorem}[Finite Alignment Theorem]
If
\[
  R_c-R_n+pL > G_i+C_i+\epsilon_i,
\]
then no modeled noncompliant action is $\epsilon_i$-optimal.  Every
$\epsilon_i$-optimal choice is policy compliant.
\end{theorem}

\begin{proof}
Let $U_c$ be compliant utility and $U_n$ the utility of any modeled
noncompliant action.  The deviation and enforcement assumptions give
\[
  U_c-U_n \ge R_c-R_n+pL-G_i-C_i.
\]
The strict-margin premise implies $U_c-U_n>\epsilon_i$.  Hence
$U_n+\epsilon_i<U_c$, contradicting the definition of an
$\epsilon_i$-optimal noncompliant choice.
\end{proof}

Strictness matters.  At equality, a noncompliant choice can remain within the
optimizer's permitted error.  The executable boundary test retains this
counterexample.

\begin{proposition}[Reward-only alignment]
Set $pL=0$.  If
\[
  R_c-R_n > G_i+C_i+\epsilon_i,
\]
then every $\epsilon_i$-optimal modeled choice is policy compliant.  Thus the
network need not punish agents when the funded advantage of ethical work alone
clears the complete deviation bound.
\end{proposition}

\begin{proposition}[Network equilibrium consequence]
If the strict-margin assumption holds for each agent and every modeled
noncompliant unilateral deviation at a profile, then no agent has an
$\epsilon_i$-optimal noncompliant deviation from that profile.  When it holds
at every reachable profile, policy compliance is a dominant action class in the
restricted game.
\end{proposition}

This proposition captures the sense in which the network can run smoothly:
routine profit seeking and policy compliance point in the same direction.
It does not remove congestion, disagreement, implementation bugs, or
unmodeled external markets.

\subsection{Coalitions and sybils}

For a coalition $K$, the analogous sufficient condition is
\[
  R_c(K)-R_n(K)+p_KL_K > G_K+C_K+\epsilon_K.
\]
An individual-agent bound does not imply this coalition bound.  Reward claims
must therefore bind a unique task or contribution nullifier.  Identity splitting
must not multiply a fixed reward.  Version 2's executable gate includes a
\texttt{taskUnclaimed} obligation; richer coalition bounds remain future work.

\section{Safety and Accounting Theorems}

The Lean development proves three complementary properties.

\begin{theorem}[Admission safety]
If the policy gate admits an action, all seven gate obligations are true.  In
particular, an admitted action is policy compliant.
\end{theorem}

\begin{theorem}[Rejected transition no-op]
If settlement rejects, payout is zero and the reward reserve is unchanged.
\end{theorem}

\begin{theorem}[Reserve conservation]
If settlement accepts, then
\[
  \text{reserve}_{post}+\text{payout}=\text{reserve}_{pre}
\]
and payout is at most the pre-settlement reserve.
\end{theorem}

These statements prevent the incentive theorem from assuming an unbacked or
unbounded reward.  They do not prove the source of reserve funds is itself
ethical; that source is another policy obligation.

\section{LLM Agents as Ethical and Profit-Seeking Workers}

An LLM agent can be trained on the network's current policy, examples of valid
work, rejected counterexamples, and the economic outcomes of accepted actions.
Its objective can combine private profit and network payout while treating hard
policy clauses as constraints:
\[
  \max_\theta\ \mathbb{E}_{a\sim\pi_\theta}
    [\text{private profit}(a)+\text{funded payout}(a)]
  \quad\text{subject to}\quad C_t(s,a,w)=1.
\]

This supports the human claim in three ways.

\begin{enumerate}
  \item \textbf{Training compatibility.}  Behavioral principles and profit
  seeking can coexist in one agent objective.  The policy gives the agent a
  precise description of acceptable work.
  \item \textbf{Economic reinforcement.}  A profitable compliant path creates
  repeated evidence that following the policy is useful rather than merely a
  restriction.
  \item \textbf{Automated coordination.}  Agents can search, negotiate, and
  complete tasks continuously, while each resulting action remains subject to
  the same publication gate as a human action.
\end{enumerate}

Present-day LLM assistants are constructive examples of the basic point: an
agent can pursue an assigned goal while following ethical and safety
constraints.  An assistant that completes useful work within those constraints
shows that the two behaviors can coexist.  This establishes feasibility at the
behavioral level.  It does not prove that every model will generalize safely to
every situation.  Version 2 does not require a particular laboratory, model
architecture, or training algorithm.  Any lab may train an agent against the
public policy and receipt interface.

The runtime shield is retained because a capable model may discover behavior
outside its training distribution or exploit a proxy reward.  The model never
grades its own evidence, chooses the active policy root, marks its own task
unclaimed, or writes the reward reserve.  Training improves how often useful
proposals pass.  The gate determines what can be committed.

\section{Tau Language Implementation}

The policy file
\texttt{tau/v2/alignment\_policy\_gate\_v2.tau} uses current Tau stream
declarations and the simple Boolean-function type.  It computes

\[
\begin{split}
\texttt{allow}={}&\texttt{policy\_root\_ok}\land
\texttt{evidence\_authenticated}\land
\texttt{action\_known}\land\\
&\texttt{action\_policy\_compliant}\land
\texttt{nonce\_fresh}\land
\texttt{task\_unclaimed}\land
\texttt{reward\_funded}.
\end{split}
\]

The replay packet contains the complete 128-row truth table.  Its first eight
rows are one all-true acceptance followed by one row that falsifies each
obligation independently.  On Tau Language 0.7.0-alpha commit
\texttt{fd137e8}, exactly one row was accepted and all 127 other rows were
rejected, matching the expected output.  The receipt binds the Tau binary and
complete packet SHA-256 hashes.  Tau remains alpha software
\cite{taulangcurrent}; this result is compatibility evidence for the bounded
kernel, not a production or mainnet claim.

The simple Boolean kernel is intentional.  Tau bitvector arithmetic is modular,
including addition, subtraction, multiplication, and division
\cite{taulangcurrent}.  Version 2 keeps exact integer accounting in the
reference settlement model and passes a verified \texttt{reward\_funded} fact
to the Tau policy.  A future bitvector settlement kernel must prove every
no-wrap precondition and preserve parity with the integer model.

\section{Lean Mechanization}

The proof project under \texttt{proofs/v2} is pinned to Lean 4.33.0 and depends
only on the Lean standard library.  The main declarations are:

\begin{itemize}
  \item \texttt{admits\_implies\_every\_obligation};
  \item \texttt{committed\_action\_is\_policy\_compliant};
  \item \texttt{rejected\_settlement\_is\_noop};
  \item \texttt{accepted\_settlement\_conserves\_reserve};
  \item \texttt{accepted\_payout\_is\_funded};
  \item \texttt{epsilon\_optimal\_choice\_is\_compliant}; and
  \item \texttt{finite\_alignment\_theorem}; and
  \item \texttt{reward\_only\_alignment\_theorem}.
\end{itemize}

Unlike Version 1, the repository includes the project configuration and exact
toolchain file.  The acceptance command is simply \texttt{lake build}.  A clean
placeholder scan is required before describing the theorem as machine checked.
The axiom audit permits only the explicitly recorded Lean standard dependencies
\texttt{propext} and \texttt{Quot.sound}; no user-added axiom is admitted.

\section{Executable Counterexample Pressure}

The exact Python reference model uses nonnegative integer atoms and rational
detection probabilities.  Its tests retain the following negative cases:

\begin{itemize}
  \item equality at the incentive boundary does not imply alignment;
  \item zero detection can make even a large nominal slash irrelevant;
  \item every individually false gate obligation rejects without an effect;
  \item zero, one-atom, exact-reserve, and maximum-u64 neighbors are checked;
  \item booleans masquerading as monetary integers and out-of-range amounts are
  rejected;
  \item a high-profit forbidden proposal cannot pass the policy shield;
  \item no admissible proposal produces a safe no-op; and
  \item tie-breaking among equal-profit admitted actions is deterministic.
\end{itemize}

An exhaustive small-domain oracle independently compares the implemented
strict-margin predicate with direct utility inequalities.  Exhaustive coverage
of this finite toy domain is evidence for that domain only.

\section{Governance and Revision}

Policy evolution is essential for a future network state.  It also creates an
authority risk.  A current policy must not silently define its own replacement
as ethical merely because the replacement is satisfiable.  Version 2 separates
an enduring constitution $H$ from revisable policy $\Phi_t$ and requires

\[
  \forall a,\quad \Phi_{t+1}(a)=1\ \Longrightarrow\ H(a)=1.
\]

The Lean project names this relation
\texttt{PreservesConstitution}.  A complete deployment needs an authenticated
update proposal, participant and voting rules, a policy-root transition,
activation time, historical receipt verification, and an emergency process
whose own authority is bounded.  Preference aggregation must publish its
identity assumptions, weights, treatment of missing information, and adversarial
manipulation bounds.  No generic scalar EETF can supply these properties by
itself.

\section{Threat Model and Nonclaims}

The principal failure modes are:

\begin{itemize}
  \item \textbf{Policy misspecification:} the formal rule omits an important
  harm or rewards a proxy;
  \item \textbf{evidence forgery:} an adapter marks false evidence as
  authenticated;
  \item \textbf{unbounded external gain:} the true benefit of deviation exceeds
  $G_i$;
  \item \textbf{governance capture:} an attacker changes the active policy or
  its identity rules;
  \item \textbf{sybil multiplication:} one contribution receives multiple
  rewards;
  \item \textbf{coalition deviation:} individually safe margins fail for a
  cooperating group;
  \item \textbf{model reward hacking:} an agent attacks the grader, evidence
  path, or reward channel; and
  \item \textbf{implementation drift:} host, Tau, Lean, documentation, and
  deployed settlement implement different relations.
\end{itemize}

Version 2 proves neither objective ethics nor the moral correctness of community
consensus.  It does not infer truth from unauthenticated sensors.  It does not
prove that an LLM's internal goals are aligned, that every future action lies in
the modeled action class, or that a private deviation bound remains valid
forever.  It supplies conditional guarantees whose assumptions can be inspected,
tested, governed, and revised.

\section{Research Program for a Future Network State}

The next steps are modular:

\begin{enumerate}
  \item define a constitutional policy vocabulary for ownership, consent,
  evidence, due process, public goods, and bounded externalities;
  \item build authenticated adapters for signatures, proofs, oracles, and task
  completion receipts;
  \item prove host--Tau--Lean parity for each policy release;
  \item measure realistic private deviation gains rather than assuming them;
  \item develop coalition, sybil, bribery, and governance-capture bounds;
  \item train agents from any capable laboratory against versioned policy and
  receipt interfaces; and
  \item run them first in shadow mode, preserving every divergence and reward
  hacking attempt as negative evidence.
\end{enumerate}

This program permits increasingly capable agents without granting them
publication authority.  As agents improve, they should become better at finding
profitable ethical work.  As policy improves, the same interface gives those
agents a clearer target.

\section{Conclusion}

Ethics and profit need not be designed as permanent adversaries.  A network can
state the behavior it will accept, fund valuable compliant work, bound the
advantage of deviation, and train automated agents to search within that space.
When the finite incentive margin is strict, policy compliance is the profitable
choice in the stated model.  When the publication gate is enforced, arbitrary
proposals cannot bypass the policy merely because an agent expects a larger
private gain.

The practical vision is unchanged from the motivating idea: ethically aligned,
profit-seeking LLM agents can automate much of a future network state's work and
help it run smoothly.  Version 2 gives that vision a narrower and stronger
foundation: explicit community policy, authenticated evidence, finite economics,
reserve-backed rewards, current Tau execution, and reproducible Lean proofs.

\bibliographystyle{plain}
\bibliography{references}

\end{document}
